\subsection{Word Sense Disambiguation}
\label{sec:word-sense-disambiguation}
\slidesrc{Semantics (long version): 31--39}

\begin{nbdef}{Word sense disambiguation (WSD)}
\textbf{Word sense disambiguation} is the task of identifying which sense (meaning) of a
word is used in a sentence, when the word has several meanings. It is an open problem of
NLP and ontology engineering.
\end{nbdef}

Consider the senses of \emph{bass} (WordNet lists eight noun senses: the lowest part of
the musical range, the lowest part in polyphonic music, an adult male singer with the
lowest voice, the lean flesh of a saltwater fish, any of various North American
freshwater fish, the member with the lowest range of a family of instruments, \ldots) and
the two sentences
\begin{itemize}[nosep]
  \item I went fishing for some sea bass.
  \item The bass line of the song is too weak.
\end{itemize}
To a human it is obvious that the first sentence uses \emph{bass} (fish) and the second
\emph{bass} (sound/instrument). Replicating this ability with an algorithm is difficult,
as shown even by the implicit equivocation between ``bass (sound)'' and ``bass (musical
instrument)'': the senses themselves are not always clearly separated.

\subsubsection{The Lesk algorithm}
\slidesrc{Semantics (long version): 35--37}

The \textbf{Lesk algorithm} (Michael Lesk, 1986, ``Automatic sense disambiguation using
machine readable dictionaries: how to tell a pine cone from an ice cream cone'') is the
classical WSD algorithm. It assumes that words in a given ``neighbourhood'' (a section of
text) tend to share a common topic.

\begin{nbdef}{Simplified Lesk}
For every sense of the word to disambiguate, count the words that occur both in the
\textbf{neighbourhood} of the word and in the \textbf{dictionary definition} (gloss) of
that sense. Choose the sense with the largest count.
\end{nbdef}

\begin{nbexample}{pine cone}
In the original version the glosses of \emph{two} ambiguous neighbouring words are
compared with each other. Context: ``pine cone''. Dictionary definitions:
\begin{center}
\small
\begin{tabular}{@{}l l@{}}
\toprule
\textbf{pine} 1 & kinds of \underline{evergreen} \underline{tree} with needle-shaped leaves \\
\textbf{pine} 2 & waste away through sorrow or illness \\
\midrule
\textbf{cone} 1 & solid body which narrows to a point \\
\textbf{cone} 2 & something of this shape whether solid or hollow \\
\textbf{cone} 3 & fruit of certain \underline{evergreen} \underline{trees} \\
\bottomrule
\end{tabular}
\end{center}
The best intersection is pine 1 $\cap$ cone 3 $=$ \{evergreen, tree(s)\}, with 2 words:
\emph{pine} is the tree and \emph{cone} is its fruit. All other pairs share no content
word.
\end{nbexample}

\begin{nbwarning}[tree vs.\ trees]
The overlap of 2 counts \emph{tree} and \emph{trees} as the same word. A literal string
comparison would find only \emph{evergreen}. Lesk needs the normalisation of
Section~\ref{sec:stemming-lemmatization} (and stop-word removal,
Section~\ref{sec:stop-words-tfidf}, otherwise \emph{of} and \emph{which} would count as
overlaps too).
\end{nbwarning}

\begin{nbexample}{simplified Lesk on ``I went fishing for some sea bass''}
After removing stop words and stemming, the context is \{go, fish, sea\}. Glosses
(with the synset members) of three senses of \emph{bass}:
\begin{itemize}[nosep]
  \item \emph{sea bass, bass}: the lean flesh of a saltwater fish of the family
        Serranidae $\to$ overlap \{sea, fish\} $= 2$;
  \item \emph{freshwater bass, bass}: any of various North American freshwater fish with
        lean flesh $\to$ overlap \{fish\} $= 1$;
  \item \emph{bass}: the lowest part of the musical range $\to$ overlap $0$.
\end{itemize}
Simplified Lesk chooses the sea bass. Note that it found \emph{sea} only thanks to the
synset member \emph{sea bass}: \emph{saltwater} and \emph{sea} do not match as strings.
\end{nbexample}

\subsubsection{Extensions and limitations}
\slidesrc{Semantics (long version): 38--39}

\begin{itemize}
  \item \textbf{Adapted Lesk} (Banerjee and Pedersen, 2002): uses WordNet, comparing
        also the glosses of related synsets (hypernyms, hyponyms, \ldots), which gives
        more words to overlap.
  \item \textbf{Limitations}: Lesk is very sensitive to the exact wording of
        definitions: the absence of a single word can change the result radically. And it
        only looks at overlaps among the glosses of the senses considered, which are
        short and do not provide enough vocabulary for fine-grained sense distinctions.
  \item Many later variants use other resources (thesauri, synonym dictionaries,
        morphological and syntactic models), and today distributional/contextual
        representations (Section~\ref{sec:distributional-semantics}) are the standard
        approach.
\end{itemize}
